\documentclass[11pt,a4paper]{article}
% Compile with XeLaTeX (fontspec and xeCJK).
\usepackage{fontspec}
\setmainfont{LibertinusSerif}[Extension=.otf,UprightFont=*-Regular,BoldFont=*-Bold,ItalicFont=*-Italic,BoldItalicFont=*-BoldItalic]
\setsansfont{LibertinusSans}[Extension=.otf,UprightFont=*-Regular,BoldFont=*-Bold,ItalicFont=*-Italic]
\usepackage{xeCJK}
\setCJKmainfont{Hiragino Mincho ProN}
\xeCJKDeclareCharClass{Default}{"00B7, "2014, "2018 -> "201F, "2026}
\usepackage[english]{babel}
\usepackage[a4paper,margin=2.5cm,top=2.3cm]{geometry}
\usepackage{graphicx}
\usepackage{xcolor}
\usepackage{titlesec}
\usepackage{fancyhdr}
\usepackage[hidelinks]{hyperref}
\definecolor{accent}{RGB}{31,73,125}
\definecolor{muted}{RGB}{93,104,119}
\hypersetup{colorlinks=true,urlcolor=accent,linkcolor=accent,pdfauthor={Romain Fernex},pdftitle={Prince Shōtoku}}
\providecommand{\tightlist}{\setlength{\itemsep}{0pt}\setlength{\parskip}{0pt}}
\providecommand{\ul}[1]{\underline{#1}}
\setcounter{secnumdepth}{0}
\titleformat{\section}{\large\bfseries\color{accent}}{}{0em}{}
\setlength{\parindent}{1.2em}
\setlength{\parskip}{0.35em}
\pagestyle{fancy}
\fancyhf{}
\fancyhead[L]{\small\color{muted}Romain Fernex}
\fancyhead[R]{\small\color{muted}\itshape Prince Shōtoku}
\fancyfoot[C]{\small\color{muted}\thepage}
\renewcommand{\headrulewidth}{0pt}
\newcommand{\articlefigure}[2]{%
  \begin{center}
  \includegraphics[width=0.55\linewidth,height=6.5cm,keepaspectratio]{#1}%
  \ifx\relax#2\relax\else\par\smallskip{\small\itshape\color{muted}#2\par}\fi
  \end{center}}

\begin{document}
\thispagestyle{empty}
\noindent{\small\sffamily\color{accent}\MakeUppercase{History of Japan}}\par\medskip
\noindent{\LARGE\bfseries Prince Shōtoku\par}
\medskip
\noindent{\color{accent}\rule{\linewidth}{0.8pt}}\par\smallskip
\noindent{\small\color{muted}Romain Fernex · 2025 · Article in English}\par\bigskip

\articlefigure{media/media/image1.jpeg}{Drawing of Shōtoku by Kikuchi Yosai}

The subject of more than one legend, Prince Shōtoku, also known as Prince Umayado, is a major figure of 7th-century Japan. Born between 572 and 574 in Nara, he was the son of Emperor Yōmei and Empress Taiza. The myth, described in the Nihon Shoki and challenged by several historians¹, holds that he was born opposite the imperial palace stables and possessed powers of clairvoyance. Made crown prince in 593 when his aunt, Empress Suiko, came to power, he is credited² with many very real transformations of Japanese society.

In politics and administration, he introduced the ``twelve-rank system of the imperial court'' in 603. Inspired by the Korean kingdom of Baekje's model, this system, although it did not truly replace hereditary succession, gave merit a greater role in officials' promotion. The prince also created the position of imperial secretary to assist the various ministries and encouraged the creation of specialised posts allocated on the basis of skills.

Internationally, in 607 he restored relations between Japan and Sui China. It was on this occasion that the expression ``Land of the Rising Sun'' is said to have appeared for the first time, in a letter addressed by Shōtoku to the emperor of China.

Finally, his religious contributions were particularly numerous. A Buddhist, he formulated seventeen injunctions in 604, linking the proper exercise of power to Buddhist values and affirming the emperor's proximity to the divine. He ordered the construction of several major temples, such as Hōryū-ji in Ikaruga (near Nara) and Shitennō-ji in Osaka. An annual ceremony in his memory, the Shōryō-e Bugaku, is still held there every 22 April. The construction of Shitennō-ji, undertaken with the support of the Kongō family --- Korean craftsmen commissioned by Shōtoku --- also gave rise to the world's oldest company, Kongō Gumi, which remains active today³.

The prince died in 622, leaving behind many children and a rich legacy blending imagined stories with historical realities.

¹ See Yoshida Kazuhiko, ``Revisioning Religion in Ancient Japan'', Japanese Journal of Religious Studies, 30/1-2,2003,p.4-5.\\
² The extent of Shōtoku's involvement in some of these advances is sometimes debated; several historians consider that Empress Suiko's role tends to be underestimated.\\
³ Although Kongō Gumi no longer exists as an independent entity, it survives as a subsidiary of the Takamatsu group, which acquired it after its bankruptcy in 2006.

\end{document}
